\documentclass{article}
\usepackage[T1]{fontenc}
\usepackage{lmodern}
\usepackage{graphicx}
\usepackage{amsmath,amssymb}
\usepackage[a4paper,margin=2cm]{geometry}
\usepackage[absolute,overlay]{textpos}
\setlength{\TPHorizModule}{1mm}
\setlength{\TPVertModule}{1mm}
\usepackage[indonesian]{babel}
\usepackage{hyperref}
\setlength{\emergencystretch}{2em}
% unit: Halaman 9

\begin{document}

% === Halaman 9 (python/native) ===

9

Mode Operasi Cipher Blok

• Mode operasi: berkaitan dengan cara blok pesan dioperasikan 

sebelum dienkripsi/dekripsi oleh fungsi E dan D.

• Ada 5 mode operasi cipher blok:

 1. Electronic Code Book (ECB)

 2. Cipher Block Chaining (CBC)

 3. Cipher Feedback (CFB)

 4.   Output Feedback (OFB)

   5.   Counter Mode

Enkripsi, E

p1

p2

...

pn

c1

c2

...

cn

Kunci K

Blok plainteks P

Blok cipherteks C

Dekripsi, D

c1

c2

...

cn

p1

p2

...

pn

Kunci K

Blok plainteks P

Blok cipherteks C

\end{document}
